\documentclass[11pt,a4paper]{article}

% ==========================================
% Pacotes de Idioma e Codificação
% ==========================================
\usepackage[portuguese]{babel}
\usepackage[utf8]{inputenc}
\usepackage[T1]{fontenc}
\usepackage{microtype} % Melhora a justificação e espaçamento do texto

% ==========================================
% Pacotes de Matemática
% ==========================================
\usepackage{amsmath}
\usepackage{amssymb}
\usepackage{amsfonts}
\usepackage{mathtools}
\usepackage{esint}      % Integrais duplos, triplos e de linha especiais
\usepackage{physics}    % Útil para diferenciais (dx, dy), derivadas, gradientes, etc.
\usepackage{cancel}     % Para simplificações de frações e termos
\usepackage{bm}         % Negrito em símbolos matemáticos
\usepackage{pgfplots}   % Criação de gráficos matemáticos
\pgfplotsset{compat=1.18}


% ==========================================
% Layout, Margens e Multicolunas
% ==========================================
\usepackage[margin=1.5cm, top=2cm, bottom=2cm]{geometry}
\usepackage{multicol}   % Permite organizar partes do texto em colunas
\usepackage{enumitem}   % Controle refinado de listas (itemize/enumerate)
\setlist{nosep}         % Remove espaços excessivos nas listas para manter compacto

% ==========================================
% Cores e Estilo Visual (Aesthetics)
% ==========================================
\usepackage{xcolor}
% Cores principais (Paleta moderna, profissional e harmoniosa)
\definecolor{primaryColor}{RGB}{20, 60, 110}   % Azul Escuro Clássico
\definecolor{secondaryColor}{RGB}{40, 110, 95} % Verde Petróleo / Teal
\definecolor{accentColor}{RGB}{180, 40, 30}    % Vermelho Terracota
\definecolor{darkGray}{RGB}{60, 60, 60}
\definecolor{lightGray}{RGB}{245, 247, 250}

% ==========================================
% Configuração de Cabeçalhos e Rodapés
% ==========================================
\usepackage{fancyhdr}
\pagestyle{fancy}
\fancyhf{}
\renewcommand{\headrulewidth}{0.4pt}
\renewcommand{\footrulewidth}{0.4pt}
\fancyhead[L]{\small\color{darkGray}\textbf{Cálculo II} --- Folha de Resumos}
\fancyhead[R]{\small\color{darkGray}Daniel Agostinho (n.º 65330)}
\fancyfoot[C]{\thepage}

% ==========================================
% Formatação de Secções e Subsecções (com linhas horizontais)
% ==========================================
\usepackage{titlesec}

% Secções (Capítulos da matéria)
\titleformat{\section}[hang]
  {\normalfont\Large\bfseries\color{primaryColor}}{Capítulo \thesection:}{0.5em}{}
\titlespacing*{\section}{0pt}{15pt}{10pt}

% Subsecções (separadas automaticamente por uma linha horizontal)
\titleformat{\subsection}
  {\normalfont\large\bfseries\color{primaryColor}}{\thesubsection}{1em}{}[\vspace{0.2em}\hrule]
\titlespacing*{\subsection}{0pt}{12pt}{6pt}

% Sub-subsecções
\titleformat{\subsubsection}
  {\normalfont\normalsize\bfseries\color{secondaryColor}}{\thesubsubsection}{1em}{}
\titlespacing*{\subsubsection}{0pt}{8pt}{4pt}

% ==========================================
% Caixas Coloridas Estruturadas (Theorems, Definitions, Rules)
% ==========================================
\usepackage[most]{tcolorbox}

% Caixa para Definições (Teal/Greenish border)
\newtcolorbox{definitionbox}[2][]{
  colback=lightGray,
  colframe=secondaryColor,
  fonttitle=\bfseries,
  title={Definição: #2},
  arc=3pt,
  left=6pt,
  right=6pt,
  top=6pt,
  bottom=6pt,
  boxrule=1.2pt,
  before=\vspace{3pt},
  after=\vspace{3pt},
  #1
}

% Caixa para Teoremas / Propriedades (Blue border)
\newtcolorbox{theorembox}[2][]{
  colback=lightGray,
  colframe=primaryColor,
  fonttitle=\bfseries,
  title={Teorema/Propriedade: #2},
  arc=3pt,
  left=6pt,
  right=6pt,
  top=6pt,
  bottom=6pt,
  boxrule=1.2pt,
  before=\vspace{3pt},
  after=\vspace{3pt},
  #1
}

% Caixa para Fórmulas Cruciais (Red border, centered math)
\newtcolorbox{formulabox}[2][]{
  colback=white,
  colframe=accentColor,
  fonttitle=\bfseries\color{white},
  colbacktitle=accentColor,
  title={Fórmula Chave: #2},
  arc=2pt,
  left=6pt,
  right=6pt,
  top=6pt,
  bottom=6pt,
  boxrule=1pt,
  halign=center,
  before=\vspace{3pt},
  after=\vspace{3pt},
  #1
}

% ==========================================
% Hiperligações (Carregar por último)
% ==========================================
\usepackage[colorlinks=true, linkcolor=primaryColor, citecolor=secondaryColor, urlcolor=secondaryColor]{hyperref}

% ==========================================
% Atalhos de Comandos Personalizados
% ==========================================
\newcommand{\R}{\mathbb{R}}
\newcommand{\N}{\mathbb{N}}
\newcommand{\Z}{\mathbb{Z}}
\newcommand{\dx}{\dd x}
\newcommand{\dy}{\dd y}
\newcommand{\dz}{\dd z}


% ==========================================
% Início do Documento
% ==========================================
\begin{document}

% ------------------------------------------
% Cabeçalho de Título Compacto
% ------------------------------------------
\begin{center}
  {\Huge\bfseries\color{primaryColor} Cálculo II}\\[0.4cm]
  {\large\bfseries Autor: Daniel Agostinho --- n.º 65330}\\[0.2cm]
  {\large\color{darkGray} Universidade Nova de Lisboa --- FCT \quad | \quad \today}\\[0.4cm]
  \rule{\linewidth}{1.5pt}
\end{center}

\vspace{0.5em}
\tableofcontents
\newpage


% ==========================================
% Capítulo 1: Equações Diferenciais Ordinárias
% ==========================================
\section{Equações Diferenciais Ordinárias}
\label{chap:edo}

\subsection{Introdução e Definições}
\begin{itemize}
  \item Seja uma função real de variável real: $\begin{aligned}[t] y: \R &\to \R \\ t &\mapsto y(t) \end{aligned}$
  \item \textbf{Taxa de variação média} (entre $t$ e $t+h$): $\frac{y(t+h) - y(t)}{h}$ (ex: velocidade, temperatura ou concentração médias).
  \item \textbf{Taxa de variação instantânea} (no ponto $t$): $y'(t) = \lim\limits_{h\to 0} \frac{y(t+h) - y(t)}{h}$ (ex: velocidade, temperatura ou concentração instantâneas).
\end{itemize}

\begin{definitionbox}{Equação Diferencial Ordinária (EDO)}
  Uma Equação Diferencial é uma equação onde a incógnita é uma função dependente de uma ou mais variáveis ($y = y(t)$) e que envolve uma ou mais derivadas desta função.
  
  Diz-se \textbf{Ordinária (EDO)} quando a função incógnita depende de apenas uma variável independente:
  \[
    F\left(t, y, y', y'', \dots, y^{(n)}\right) = 0
  \]
  \textbf{Ordem:} É a ordem da maior derivada envolvida na equação.\\
  \textbf{Solução:} Função que satisfaz a EDO num certo aberto $I$ (pode ser explícita ou implícita).\\
  \textbf{Curva Integral:} O gráfico de uma solução da EDO. Uma solução geral produz uma família de curvas integrais.
  \begin{center}
    \begin{tikzpicture}[scale=1]
      \begin{axis}[
        domain=-2:2,
        axis lines=middle,
        xlabel={$x$},
        ylabel={$y$},
        title={\small Família de Curvas Integrais: $y' = 2x \implies y = x^2 + C$},
        width=5.5cm, height=5cm,
        tick label style={font=\tiny},
        label style={font=\tiny},
        title style={font=\tiny\bfseries\color{primaryColor}},
        axis line style={gray, -stealth},
        ymin=-3, ymax=3
      ]
        \addplot[thick, blue!90] {x^2 - 2};
        \addplot[thick, blue!90] {x^2 - 1};
        \addplot[thick, blue!90] {x^2};
        \addplot[thick, blue!90] {x^2 + 1};
      \end{axis}
    \end{tikzpicture}
  \end{center}
\end{definitionbox}

\begin{definitionbox}{Classificação das Soluções}
  \begin{itemize}
    \item \textbf{Solução Geral:} Expressão com constantes arbitrárias que engloba todas as soluções (exceto, no máximo, um número finito). Para ordem $n$, apresenta $n$ constantes arbitrárias.
    \item \textbf{Solução Particular:} Solução concreta obtida a partir da solução geral especificando valores para as constantes.
    \item \textbf{Solução Singular:} Solução que não pode ser obtida a partir da geral por particularização das constantes.
    \item \textbf{Solução Estável ou de Equilíbrio:} Solução constante da EDO (i.e., $y(t) = c \implies y' = 0$).
  \end{itemize}
\end{definitionbox}

\newpage

\subsection{Problema de Valor Inicial (PVI)}
A solução geral de uma EDO de ordem $n$ apresenta $n$ constantes arbitrárias. Para concretizar uma solução particular, define-se um PVI:
\[
  \begin{cases}
    y' = f(x, y) \\
    y(x_0) = y_0
  \end{cases}
\]
Não existe um método geral para resolver EDOs de 1.ª ordem. Usam-se técnicas específicas para certos tipos.

\begin{theorembox}{Teorema de Existência e Unicidade (Picard-Lindelöf)}
  Se $f(x, y)$ e $\frac{\partial f}{\partial y}$ forem contínuas num retângulo aberto contendo $(x_0, y_0)$, então o PVI tem uma única solução $y(x)$ definida numa vizinhança de $x_0$.
\end{theorembox}

\subsection{Resolução de EDOs de 1.ª Ordem}

\begin{theorembox}{EDO Linear de 1.ª Ordem}
  EDO da forma $y' + g(x) y = f(x)$ onde $g$ e $f$ são contínuas. \\ Resolve-se pelo \textbf{Método do Fator Integrante}:
  \begin{enumerate}
    \item Fator integrante: $\mu(x) = e^{G(x)} = e^{\int g(x) \dx}$ onde $G(x) = \int g(x) \dx$.
    \item Multiplica-se a EDO por $\mu(x)$: $(\mu(x) y)' = \mu(x) f(x)$.
    \item Solução geral: $y(x) = \frac{1}{\mu(x)} \left[ \int \mu(x) f(x) \dx + C \right]$.
  \end{enumerate}
  \textbf{Exemplo Prático (Slide 1):} $t y' + (1+t)y = 2$ para $t > 0$.
  \begin{enumerate}
    \item Forma padrão: $y' + \frac{1+t}{t}y = \frac{2}{t} \implies g(t) = 1 + \frac{1}{t}$ e $f(t) = \frac{2}{t}$.
    \item Fator integrante: $\mu(t) = e^{\int \left(1 + \frac{1}{t}\right) \dd t} = e^{t + \ln t} = t e^t$.
    \item Multiplicar por $\mu(t)$: $(t e^t y)' = 2 e^t$.
    \item Integrando: $t e^t y = 2 e^t + C \implies y(t) = \frac{2}{t} + \frac{C}{t} e^{-t}$.
  \end{enumerate}
\end{theorembox}

\begin{theorembox}{EDO de Variáveis Separáveis}
  EDO da forma $h(y) \frac{\dd y}{\dd x} = g(x)$, onde $g$ e $h$ são contínuas: $\int h(y) \dd y = \int g(x) \dx + C \implies H(y) = G(x) + C$.
  
  \textbf{Aviso de Perda de Soluções:} 
  
  No processo de divisão para separar variáveis, assumindo $h(y) \neq 0$, podem perder-se soluções constantes (singulares).
  
  \textbf{Exemplo Prático (Slide 1):} Resolver a EDO $y' - 3y = 0$:
  \begin{enumerate}
    \item Separar as variáveis: $y' = 3y \implies \frac{1}{y} \dd y = 3 \dd x$.
    \item Integrar ambos os membros: $\int \frac{1}{y} \dd y = \int 3 \dd x \implies \ln|y| = 3x + C_1$.
    \item Resolver em ordem a $y$: $|y| = e^{C_1} e^{3x} \implies y(x) = \pm e^{C_1} e^{3x}$.
    \item Definir a constante geral: $y(x) = C_0 e^{3x}$ (com $C_0 \neq 0$).
    \item Análise da Solução Perdida: A função constante $y(x) = 0$ é solução da EDO original, mas foi perdida ao assumirmos $y \neq 0$. Ela é uma \textbf{solução singular}. Se alargarmos a constante para $C_0 \in \R$, a solução singular fica englobada na expressão geral.
  \end{enumerate}
\end{theorembox}

\subsection{Modelos de Crescimento, Decaimento e Temperatura}

\begin{formulabox}{Modelos Exponenciais}
  Taxa de variação proporcional à quantidade atual: $\frac{\dd y}{\dd x} = k y$ ($k \neq 0$). Solução geral: $y(x) = C e^{kx}$.
  \begin{itemize}
    \item $k > 0$: Crescimento exponencial (ex: juros, populações, vírus).
    \item $k < 0$: Decaimento exponencial (ex: decaimento radioativo).
    \item \textbf{Meia-vida ($T_{1/2}$):} Tempo necessário para a quantidade decair para metade da inicial:
    \begin{align*}
      y(T_{1/2}) = \frac{1}{2}y(0) &\iff y(0)e^{k T_{1/2}} = \frac{1}{2}y(0) \\
      &\iff e^{k T_{1/2}} = \frac{1}{2} \\
      &\iff k T_{1/2} = \ln\left(\frac{1}{2}\right) = -\ln 2 \\
      &\iff k = - \frac{\ln 2}{T_{1/2}}
    \end{align*}
  \end{itemize}
  \textbf{Exemplo 1 (Plutónio-239):} Sabendo que $T_{1/2} = 24100$ anos, quanto tempo demoram $10$\,gr a decair para $1$\,gr?
  \[
    k = - \frac{\ln 2}{24100} \implies 1 = 10 e^{kt} \implies \ln(0.1) = kt \implies t = \frac{\ln 0.1}{-\frac{\ln 2}{24100}} = 24100 \frac{\ln 10}{\ln 2} \approx 80058 \text{ anos}.
  \]
  \textbf{Exemplo 2 (Carbono-14):} Fóssil retém $15\%$ do C-14 atual. Meia-vida $T_{1/2} = 5730$ anos. Idade da amostra:
  \[
    k = - \frac{\ln 2}{5730} \implies 0.15 = 1 \cdot e^{kt} \implies t = \frac{\ln 0.15}{k} = 5730 \frac{\ln(1/0.15)}{\ln 2} \approx 15683 \text{ anos}.
  \]
\end{formulabox}

\begin{formulabox}{Lei da Variação de Temperatura de Newton}
  A taxa de variação da temperatura $T(t)$ de um objeto é proporcional à diferença entre a sua temperatura e a temperatura constante do ambiente $T_a$:
  \[
    \frac{\dd T}{\dd t} = k(T - T_a) \implies \int \frac{\dd T}{T-T_a} = \int k \dd t \implies T(t) = T_a + C e^{kt}
  \]
  \textbf{Exemplo 3 (Copo de Água):} Copo a $95^\circ$ em sala a $21^\circ$. Ao fim de $1$\,min, o copo está a $85^\circ$. Tempo para atingir $51^\circ$:
  \begin{enumerate}
    \item $T(t) = 21 + C e^{kt} \implies T(0) = 95 \implies C = 95 - 21 = 74 \implies T(t) = 21 + 74 e^{kt}$.
    \item $T(1) = 85 \implies 85 = 21 + 74 e^{k} \implies e^k = \frac{64}{74} = \frac{32}{37} \implies k = \ln(32/37) \approx -0.145$.
    \item $T(t) = 51 \implies 51 = 21 + 74 e^{kt} \implies e^{kt} = \frac{30}{74} = \frac{15}{37} \implies t = \frac{\ln(15/37)}{\ln(32/37)} \approx 6.2 \text{ minutos}$.
  \end{enumerate}
\end{formulabox}

\newpage

\subsection{Campos de Direções e Métodos Numéricos}
\begin{itemize}
  \item \textbf{Campo de Direções:} Representação gráfica em malha bidimensional dos declives (inclinações) das curvas integrais dados por $y' = f(x, y)$ nos pontos $(x, y)$. Permite visualizar qualitativamente as soluções. Exemplo: $y' = y - x$.
\end{itemize}

\begin{center}
  \begin{minipage}{0.5\textwidth}
    \centering
    \begin{tikzpicture}[scale=0.9]
      \begin{axis}[
        view={0}{90},
        domain=-2:2,
        y domain=-2:2,
        samples=9,
        axis lines=middle,
        xlabel={$x$},
        ylabel={$y$},
        title={\small Campo de Direções: $y' = y - x$},
        width=6.5cm, height=6.5cm,
        tick label style={font=\tiny},
        label style={font=\small},
        title style={font=\small\bfseries\color{primaryColor}},
        axis line style={gray, -stealth}
      ]
        \addplot3[
          blue!80!black,
          quiver={
            u={1/sqrt(1 + (y-x)^2)},
            v={(y-x)/sqrt(1 + (y-x)^2)},
            scale arrows=0.2
          },
          -stealth,
          samples=9
        ] {0};
      \end{axis}
    \end{tikzpicture}
  \end{minipage}%
  \begin{minipage}{0.5\textwidth}
    \centering
    \small
    \textbf{Tabela de Declives ($y' = y - x$):}
    \vspace{0.5em}
    
    \begin{tabular}{cc|c}
      \hline
      $x$ & $y$ & $y' = y - x$ \\
      \hline
      $0$ & $0$ & $0$ \\
      $1$ & $0$ & $-1$ \\
      $-1$ & $0$ & $1$ \\
      $0$ & $1$ & $1$ \\
      $0$ & $-1$ & $-1$ \\
      $1$ & $1$ & $0$ \\
      $2$ & $1$ & $-1$ \\
      $1$ & $2$ & $1$ \\
      \hline
    \end{tabular}
  \end{minipage}
\end{center}

\begin{theorembox}{Método de Euler}
  Método numérico iterativo para aproximar a solução de um PVI $y' = f(x, y), y(x_0) = y_0$ em pontos distanciados entre si uma quantidade fixa $h > 0$ (passo):
  \begin{align*}
    x_{n+1} &= x_n + h \\
    y_{n+1} &= y_n + h f(x_n, y_n)
  \end{align*}
  \textbf{Nota:} A qualidade da aproximação melhora progressivamente à medida que o tamanho do passo $h$ é reduzido.
  
  \begin{center}
    \begin{tikzpicture}[scale=1.0]
      % Axes
      \draw[-stealth, thick, gray] (-0.5,0) -- (4,0) node[right, black] {\tiny $x$};
      \draw[-stealth, thick, gray] (0,-0.5) -- (0,3.5) node[above, black] {\tiny $y$};
      
      % Exact curve (smooth curve)
      \draw[blue, thick] (0.5,0.8) to[out=20, in=250] (2,2.0) to[out=70, in=200] (3.5,3.2);
      \node[blue, right] at (3.5,3.2) {\tiny Exata $y(x)$};
      
      % Euler points and lines
      \coordinate (P0) at (0.5,0.8);
      \coordinate (P1) at (1.5,1.1);
      \coordinate (P2) at (2.5,1.8);
      \coordinate (P3) at (3.5,2.9);
      
      \draw[red, thick, mark=*] (P0) -- (P1) -- (P2) -- (P3);
      
      % Points labels
      \fill[red] (P0) circle (1.5pt) node[below right, black] {\tiny $(x_0, y_0)$};
      \fill[red] (P1) circle (1.5pt) node[below right, black] {\tiny $(x_1, y_1)$};
      \fill[red] (P2) circle (1.5pt) node[below right, black] {\tiny $(x_2, y_2)$};
      \fill[red] (P3) circle (1.5pt) node[below right, black] {\tiny $(x_3, y_3)$};
      
      % Step size labels (dashed lines)
      \draw[dashed, gray] (0.5,0) -- (0.5,0.8);
      \draw[dashed, gray] (1.5,0) -- (1.5,1.1);
      \draw[dashed, gray] (2.5,0) -- (2.5,1.8);
      \draw[dashed, gray] (3.5,0) -- (3.5,2.9);
      
      \node[below] at (0.5,0) {\tiny $x_0$};
      \node[below] at (1.5,0) {\tiny $x_1$};
      \node[below] at (2.5,0) {\tiny $x_2$};
      \node[below] at (3.5,0) {\tiny $x_3$};
      
      % Horizontal arrows for step h
      \draw[cyan, -stealth] (0.5, -0.3) -- (1.5, -0.3) node[midway, below] {\tiny $h$};
      \draw[cyan, -stealth] (1.5, -0.3) -- (2.5, -0.3) node[midway, below] {\tiny $h$};
      
      \node[red, above left] at (2.5,1.8) {\tiny Euler};
    \end{tikzpicture}
  \end{center}
\end{theorembox}

\newpage


% ==========================================
% Capítulo 2: Cónicas e Quádricas
% ==========================================
\section{Cónicas e Quádricas}
\label{chap:conicas_quadricas}

\subsection{Cónicas}

\begin{definitionbox}{Secção Cónica}
  Uma \textbf{secção cónica} é a curva obtida através da interseção de um plano bidimensional com um cone circular reto duplo. 
  Analiticamente, qualquer cónica no plano $xOy$ é definida por uma equação geral de segundo grau:
  \begin{equation}
    Ax^2 + By^2 + Cxy + Dx + Ey + F = 0 \label{eq:conica_geral}
  \end{equation}
  \textbf{Classificação por Termos:}
  \begin{itemize}
    \item \textbf{Figuras Normais ($C = 0$):} Os eixos da cónica são paralelos aos eixos coordenados. O centro ou vértice está na origem se $D = E = 0$.
    \item \textbf{Figuras Rodadas ($C \neq 0$):} A presença do termo misto $Cxy$ indica uma rotação dos eixos da cónica em relação ao sistema $xOy$.
  \end{itemize}
\end{definitionbox}

\subsubsection{Elipse}
\begin{minipage}{0.58\textwidth}
  \begin{definitionbox}{Elipse}
    Conjunto dos pontos do plano cuja soma das distâncias a dois pontos fixos $F_1$ e $F_2$ (focos), distanciados de $2c$, é constante e igual a $2a$, com $a > c$.
    
    \textbf{Relação Fundamental:}
    \begin{equation}
      a^2 = b^2 + c^2 \label{eq:elipse_rel}
    \end{equation}
    Onde $a$ e $b$ são os semi-eixos maior e menor.
    
    \textbf{Equação Normal (Centro na Origem $O(0,0)$):}
    \begin{equation}
      \frac{x^2}{a^2} + \frac{y^2}{b^2} = 1 \quad (\text{Eixo maior ao longo de } OX)
    \end{equation}
    \textbf{Equação Geral (Centro em $C(h,k)$):}
    \begin{equation}
      \frac{(x-h)^2}{a^2} + \frac{(y-k)^2}{b^2} = 1
    \end{equation}
    \textit{Nota:} Se $b > a$, o eixo maior é vertical ($OY$), os focos são $(h, k \pm c)$, e a relação é $b^2 = a^2 + c^2$.
  \end{definitionbox}
\end{minipage}
\hfill
\begin{minipage}{0.38\textwidth}
  \centering
  \begin{tikzpicture}[scale=1.1]
    % Draw axes
    \draw[-stealth, gray] (-2.3,0) -- (2.3,0) node[right, black] {\tiny $x$};
    \draw[-stealth, gray] (0,-1.5) -- (0,1.5) node[above, black] {\tiny $y$};
    
    % Draw ellipse: a=1.8, b=1.1 => c = sqrt(3.24 - 1.21) = sqrt(2.03) \approx 1.42
    \draw[thick, primaryColor] (0,0) ellipse (1.8cm and 1.1cm);
    
    % Foci
    \coordinate (F1) at (-1.42,0);
    \coordinate (F2) at (1.42,0);
    \fill[accentColor] (F1) circle (1.5pt) node[below] {\tiny $F_1$};
    \fill[accentColor] (F2) circle (1.5pt) node[below] {\tiny $F_2$};
    
    % Vertices
    \fill (1.8,0) circle (1pt) node[above right] {\tiny $V_2(a,0)$};
    \fill (-1.8,0) circle (1pt) node[above left] {\tiny $V_1$};
    \fill (0,1.1) circle (1pt) node[above left] {\tiny $V_3(0,b)$};
    \fill (0,-1.1) circle (1pt) node[below left] {\tiny $V_4$};
    
    % Right triangle representing a^2 = b^2 + c^2
    \draw[dashed, darkGray] (0,0) -- (0,1.1) node[midway, left] {\tiny $b$};
    \draw[dashed, darkGray] (0,0) -- (1.42,0) node[midway, below] {\tiny $c$};
    \draw[thick, secondaryColor] (0,1.1) -- (1.42,0) node[midway, above right] {\tiny $a$};
    
    % Center
    \node[below left] at (0,0) {\tiny $O$};
  \end{tikzpicture}
  \\[0.2cm]
  \small\textbf{Figura 2.1:} Elementos da Elipse.
\end{minipage}

\subsubsection{Hipérbole}
\begin{minipage}{0.58\textwidth}
  \begin{definitionbox}{Hipérbole}
    Conjunto dos pontos do plano cuja diferença absoluta das distâncias a dois pontos fixos $F_1$ e $F_2$ (focos), distanciados de $2c$, é constante e igual a $2a$, com $a < c$.
    
    \textbf{Relação Fundamental:}
    \begin{equation}
      c^2 = a^2 + b^2 \label{eq:hiperbole_rel}
    \end{equation}
    Onde $a$ é o semi-eixo real e $b$ o semi-eixo imaginário.
    
    \textbf{Equação Normal (Centro na Origem $O(0,0)$):}
    \begin{equation}
      \frac{x^2}{a^2} - \frac{y^2}{b^2} = 1 \quad (\text{Abertura horizontal, eixo real } OX)
    \end{equation}
    \textbf{Assíntotas:} $y = \pm \frac{b}{a} x$.
    
    \textbf{Equação Geral (Centro em $C(h,k)$):}
    \begin{equation}
      \frac{(x-h)^2}{a^2} - \frac{(y-k)^2}{b^2} = 1
    \end{equation}
    \textit{Nota:} Se a abertura for vertical, a equação normal é $\frac{y^2}{b^2} - \frac{x^2}{a^2} = 1$, com assíntotas $y = \pm \frac{b}{a} x$ e focos em $(0, \pm c)$.
  \end{definitionbox}
\end{minipage}
\hfill
\begin{minipage}{0.38\textwidth}
  \centering
  \begin{tikzpicture}[scale=1.0]
    % Draw axes
    \draw[-stealth, gray] (-2.5,0) -- (2.5,0) node[right, black] {\tiny $x$};
    \draw[-stealth, gray] (0,-1.8) -- (0,1.8) node[above, black] {\tiny $y$};
    
    % Asymptotes: y = +/- (b/a)x. Let a=1.1, b=0.9 => c = sqrt(1.21+0.81) = sqrt(2.02) \approx 1.42
    \draw[dashed, gray] (-2.0,-1.63) -- (2.0,1.63);
    \draw[dashed, gray] (-2.0,1.63) -- (2.0,-1.63);
    
    % Auxiliary rectangle: 2a x 2b
    \draw[dotted, darkGray] (-1.1,-0.9) rectangle (1.1,0.9);
    \fill (1.1,0) circle (1pt) node[below right] {\tiny $V_2$};
    \fill (-1.1,0) circle (1pt) node[below left] {\tiny $V_1$};
    
    % Foci
    \coordinate (F1) at (-1.42,0);
    \coordinate (F2) at (1.42,0);
    \fill[accentColor] (F1) circle (1.5pt) node[below] {\tiny $F_1$};
    \fill[accentColor] (F2) circle (1.5pt) node[below] {\tiny $F_2$};
    
    % Hyperbola curves
    \draw[thick, primaryColor, domain=-50:50, samples=100] 
      plot ({1.1/cos(\x)}, {0.9*tan(\x)});
    \draw[thick, primaryColor, domain=-50:50, samples=100] 
      plot ({-1.1/cos(\x)}, {0.9*tan(\x)});
    
    % Triangle for c^2 = a^2 + b^2
    \draw[thick, secondaryColor] (0,0) -- (1.1,0.9) node[midway, above left] {\tiny $c$};
    \draw[dashed, darkGray] (0,0) -- (1.1,0) node[midway, below] {\tiny $a$};
    \draw[dashed, darkGray] (1.1,0) -- (1.1,0.9) node[midway, right] {\tiny $b$};
  \end{tikzpicture}
  \\[0.2cm]
  \small\textbf{Figura 2.2:} Elementos da Hipérbole.
\end{minipage}

\subsubsection{Parábola}
\begin{minipage}{0.58\textwidth}
  \begin{definitionbox}{Parábola}
    Conjunto dos pontos do plano equidistantes de um ponto fixo $F$ (foco) e de uma reta fixa $d$ (diretriz) que não contém $F$: $d(P,F) = d(P,d)$.
    
    \textbf{Foco e Diretriz:} A distância do vértice ao foco é $p = \frac{1}{4a}$.
    
    \textbf{Equação Normal (Vértice na Origem $V(0,0)$, Eixo Vertical):}
    \begin{equation}
      y = ax^2 \quad \left(F\left(0, \frac{1}{4a}\right), \quad d: y = -\frac{1}{4a}\right)
    \end{equation}
    \begin{itemize}
      \item $a > 0$: Concavidade voltada para cima.
      \item $a < 0$: Concavidade voltada para baixo.
    \end{itemize}
    \textbf{Equações Gerais (Vértice em $V(h,k)$):}
    \begin{align}
      \text{Eixo Vertical:} \quad & y - k = a(x-h)^2 \\
      \text{Eixo Horizontal:} \quad & x - h = a(y-k)^2
    \end{align}
  \end{definitionbox}
\end{minipage}
\hfill
\begin{minipage}{0.38\textwidth}
  \centering
  \begin{tikzpicture}[scale=1.1]
    % Draw axes
    \draw[-stealth, gray] (-2.0,0) -- (2.0,0) node[right, black] {\tiny $x$};
    \draw[-stealth, gray] (0,-0.8) -- (0,2.2) node[above, black] {\tiny $y$};
    
    % Parabola: y = a x^2. Let a = 0.5 => p = 0.5
    \draw[thick, primaryColor, domain=-1.8:1.8, samples=50] plot (\x, {0.5*\x*\x});
    
    % Focus
    \coordinate (F) at (0,0.5);
    \fill[accentColor] (F) circle (1.5pt) node[left] {\tiny $F(0,p)$};
    
    % Directrix line: y = -0.5
    \draw[thick, dashed, secondaryColor] (-1.9,-0.5) -- (1.9,-0.5) node[right] {\tiny $d$};
    
    % Point P on parabola: P = (1.2, 0.5 * 1.2^2) = (1.2, 0.72)
    \coordinate (P) at (1.2,0.72);
    \fill (P) circle (1.5pt) node[above right] {\tiny $P(x,y)$};
    
    % Distances to focus and directrix
    \draw[dotted, darkGray, thick] (P) -- (F);
    \draw[dotted, darkGray, thick] (P) -- (1.2,-0.5);
    \fill (1.2,-0.5) circle (1pt);
  \end{tikzpicture}
  \\[0.2cm]
  \small\textbf{Figura 2.3:} Elementos da Parábola.
\end{minipage}

\newpage

\subsection{Espaço Tridimensional e Quádricas}

\subsubsection{Espaço Tridimensional \texorpdfstring{$\mathbb{R}^3$}{R\textasciicircum 3} e Secções}
No espaço $\mathbb{R}^3$, pontos são representados por triplos ordenados $(x,y,z)$.
\begin{itemize}
  \item \textbf{Planos Gerais:} $Ax + By + Cz + D = 0$.
  \item \textbf{Planos de Secção:} Planos paralelos aos planos coordenados são obtidos fixando variáveis: $x = a$ (paralelo a $yOz$), $y = b$ (paralelo a $xOz$) e $z = c$ (paralelo a $xOy$).
\end{itemize}

\begin{center}
  \begin{tikzpicture}[x={(-0.5cm,-0.5cm)}, y={(1.0cm,0cm)}, z={(0cm,1.0cm)}, scale=1.1]
    % --- Plano x = a ---
    \begin{scope}[shift={(0,0,0)}]
      % Axes
      \draw[-stealth, gray] (0,0,0) -- (0,2.0,0) node[right, black] {\tiny $y$};
      \draw[-stealth, gray] (0,0,0) -- (0,0,2.0) node[above, black] {\tiny $z$};
      \draw[-stealth, gray] (0,0,0) -- (2.0,0,0) node[below left, black] {\tiny $x$};
      
      % Plane x = 1.0 (parallel to yOz)
      \filldraw[fill=primaryColor!15, draw=primaryColor!80, opacity=0.75, thick] 
        (1.0, 0, 0) -- (1.0, 1.8, 0) -- (1.0, 1.8, 1.8) -- (1.0, 0, 1.8) -- cycle;
      
      \fill[black] (1.0,0,0) circle (1.5pt) node[below left] {\tiny $(a,0,0)$};
      \node[above] at (1.0,0.9,1.8) {\small Plano $x = a$};
    \end{scope}

    % --- Plano y = b ---
    \begin{scope}[shift={(0,3.5,0)}]
      % Axes
      \draw[-stealth, gray] (0,0,0) -- (0,2.0,0) node[right, black] {\tiny $y$};
      \draw[-stealth, gray] (0,0,0) -- (0,0,2.0) node[above, black] {\tiny $z$};
      \draw[-stealth, gray] (0,0,0) -- (2.0,0,0) node[below left, black] {\tiny $x$};
      
      % Plane y = 1.0 (parallel to xOz)
      \filldraw[fill=secondaryColor!15, draw=secondaryColor!80, opacity=0.75, thick] 
        (0, 1.0, 0) -- (1.8, 1.0, 0) -- (1.8, 1.0, 1.8) -- (0, 1.0, 1.8) -- cycle;
      
      \fill[black] (0,1.0,0) circle (1.5pt) node[below right] {\tiny $(0,b,0)$};
      \node[above] at (0.9,1.0,1.8) {\small Plano $y = b$};
    \end{scope}

    % --- Plano z = c ---
    \begin{scope}[shift={(0,7.0,0)}]
      % Axes
      \draw[-stealth, gray] (0,0,0) -- (0,2.0,0) node[right, black] {\tiny $y$};
      \draw[-stealth, gray] (0,0,0) -- (0,0,2.0) node[above, black] {\tiny $z$};
      \draw[-stealth, gray] (0,0,0) -- (2.0,0,0) node[below left, black] {\tiny $x$};
      
      % Plane z = 1.0 (parallel to xOy)
      \filldraw[fill=accentColor!10, draw=accentColor!80, opacity=0.75, thick] 
        (0, 0, 1.0) -- (1.8, 0, 1.0) -- (1.8, 1.8, 1.0) -- (0, 1.8, 1.0) -- cycle;
      
      \fill[black] (0,0,1.0) circle (1.5pt) node[left] {\tiny $(0,0,c)$};
      \node[above] at (0.9,0.9,1.0) {\small Plano $z = c$};
    \end{scope}
  \end{tikzpicture}
\end{center}

\subsubsection{Superfícies Quádricas}
Uma \textbf{superfície quádrica} é definida analiticamente por uma equação tridimensional de segundo grau:
\begin{equation}
  Ax^2 + By^2 + Cz^2 + Dxy + Exz + Fyz + Gx + Hy + Iz + J = 0 \label{eq:quadrica_geral}
\end{equation}
\textbf{Formas Normais ($D = E = F = 0$):} A superfície tem eixos paralelos aos eixos cartesianos e centro ou vértice na origem.

\begin{formulabox}[top=2pt, bottom=2pt, left=4pt, right=4pt]{1. Superfície Esférica (Esfera)}
  \begin{minipage}{0.72\textwidth}
    \textbf{Forma Normal:} $x^2 + y^2 + z^2 = r^2$ \\
    \textbf{Translada em $C(h,j,l)$:} $(x-h)^2 + (y-j)^2 + (z-l)^2 = r^2$ \\
    \textbf{Características:} Centro $C(0,0,0)$ e raio $r$. As secções por planos $x=k$, $y=k$ ou $z=k$ são circunferências de raio $\sqrt{r^2-k^2}$ (para $|k| < r$).
  \end{minipage}%
  \hfill
  \begin{minipage}{0.24\textwidth}
    \centering
    \begin{tikzpicture}[scale=0.52]
      \draw[-stealth, gray] (-1.5,0) -- (1.8,0) node[right, black] {\tiny $y$};
      \draw[-stealth, gray] (0,-1.5) -- (0,1.8) node[above, black] {\tiny $z$};
      \draw[-stealth, gray] (0.8,-0.8) -- (-0.8,0.8) node[above left, black] {\tiny $x$};
      
      \draw[thick, primaryColor] (0,0) circle (1.1);
      \draw[dashed, secondaryColor] (1.1,0) arc (0:180:1.1 and 0.35);
      \draw[thick, secondaryColor] (-1.1,0) arc (180:360:1.1 and 0.35);
      \draw[dashed, secondaryColor] (0,1.1) arc (90:270:0.35 and 1.1);
      \draw[thick, secondaryColor] (0,-1.1) arc (-90:90:0.35 and 1.1);
    \end{tikzpicture}
  \end{minipage}
\end{formulabox}

\begin{formulabox}[top=2pt, bottom=2pt, left=4pt, right=4pt]{2. Elipsóide}
  \begin{minipage}{0.72\textwidth}
    \textbf{Forma Normal:} $\dfrac{x^2}{a^2} + \dfrac{y^2}{b^2} + \dfrac{z^2}{c^2} = 1$ \\
    \textbf{Translada em $C(h,j,l)$:} $\dfrac{(x-h)^2}{a^2} + \dfrac{(y-j)^2}{b^2} + \dfrac{(z-l)^2}{c^2} = 1$ \\
    \textbf{Características:} Simétrica em relação a todos os planos coordenados. As secções por planos paralelos aos eixos ($x=k$, $y=k$, $z=k$) são elipses ou circunferências.
  \end{minipage}%
  \hfill
  \begin{minipage}{0.24\textwidth}
    \centering
    \begin{tikzpicture}[scale=0.52]
      \draw[-stealth, gray] (-1.7,0) -- (2.0,0) node[right, black] {\tiny $y$};
      \draw[-stealth, gray] (0,-1.3) -- (0,1.6) node[above, black] {\tiny $z$};
      \draw[-stealth, gray] (0.8,-0.8) -- (-0.8,0.8) node[above left, black] {\tiny $x$};
      
      \draw[thick, primaryColor] (0,0) ellipse (1.4 and 0.9);
      \draw[dashed, secondaryColor] (1.4,0) arc (0:180:1.4 and 0.3);
      \draw[thick, secondaryColor] (-1.4,0) arc (180:360:1.4 and 0.3);
    \end{tikzpicture}
  \end{minipage}
\end{formulabox}

\begin{formulabox}[top=2pt, bottom=2pt, left=4pt, right=4pt]{3. Parabolóide Elíptico}
  \begin{minipage}{0.72\textwidth}
    \textbf{Forma Normal (eixo $OZ$):} $z = \dfrac{x^2}{a^2} + \dfrac{y^2}{b^2}$ \\
    \textbf{Translada em $V(h,j,k_0)$:} $z-k_0 = \dfrac{(x-h)^2}{a^2} + \dfrac{(y-j)^2}{b^2}$ \\
    \textbf{Características:} O vértice é $V(0,0,0)$. 
    As secções $x=k$ e $y=k$ são parábolas. As secções $z=k$ são elipses para $k>0$ (nulas se $k<0$). \\
    \textit{Permutações:} Se a variável linear for $x$ (eixo $OX$) ou $y$ (eixo $OY$), o parabolóide abre-se ao longo desse eixo.
  \end{minipage}%
  \hfill
  \begin{minipage}{0.24\textwidth}
    \centering
    \begin{tikzpicture}[scale=0.52]
      \draw[-stealth, gray] (-1.5,0) -- (1.8,0) node[right, black] {\tiny $y$};
      \draw[-stealth, gray] (0,-0.2) -- (0,2.0) node[above, black] {\tiny $z$};
      \draw[-stealth, gray] (0.8,-0.8) -- (-0.8,0.8) node[above left, black] {\tiny $x$};
      
      \draw[thick, primaryColor] (-1.0,1.4) to[out=-80, in=180] (0,0) to[out=0, in=-100] (1.0,1.4);
      \draw[thick, secondaryColor] (0,1.4) ellipse (1.0 and 0.3);
      \draw[dashed, gray] (-0.5,0.5) to[out=-60, in=180] (0,0.1) to[out=0, in=-120] (0.5,0.5);
    \end{tikzpicture}
  \end{minipage}
\end{formulabox}

\begin{formulabox}[top=2pt, bottom=2pt, left=4pt, right=4pt]{4. Parabolóide Hiperbólico (Sela)}
  \begin{minipage}{0.72\textwidth}
    \textbf{Forma Normal (eixo $OZ$):} $z = \dfrac{x^2}{a^2} - \dfrac{y^2}{b^2}$ \\
    \textbf{Translada em $V(h,j,k_0)$:} $z-k_0 = \dfrac{(x-h)^2}{a^2} - \dfrac{(y-j)^2}{b^2}$ \\
    \textbf{Características:} O ponto de sela (máximo numa direção e mínimo noutra) é $V(0,0,0)$. \\
    As secções $x=k$ e $y=k$ são parábolas (uma com concavidade para cima, outra para baixo). As secções $z=k$ são hipérboles.
  \end{minipage}%
  \hfill
  \begin{minipage}{0.24\textwidth}
    \centering
    \begin{tikzpicture}[scale=0.52]
      \draw[-stealth, gray] (-1.5,0) -- (1.8,0) node[right, black] {\tiny $y$};
      \draw[-stealth, gray] (0,-1.4) -- (0,1.8) node[above, black] {\tiny $z$};
      \draw[-stealth, gray] (0.8,-0.8) -- (-0.8,0.8) node[above left, black] {\tiny $x$};
      
      \draw[thick, primaryColor] (-1.2,-0.6) to[out=60, in=180] (0,0.4) to[out=0, in=120] (1.2,-0.6);
      \draw[thick, primaryColor] (-0.8,0.6) to[out=-45, in=180] (0,-0.4) to[out=0, in=-135] (0.8,0.6);
      \draw[dashed, secondaryColor] (-1.2,-0.6) -- (-0.8,0.6);
      \draw[thick, secondaryColor] (1.2,-0.6) -- (0.8,0.6);
    \end{tikzpicture}
  \end{minipage}
\end{formulabox}

\begin{formulabox}[top=2pt, bottom=2pt, left=4pt, right=4pt]{5. Cone Elíptico}
  \enlargethispage{1.8cm}
  \begin{minipage}{0.72\textwidth}
    \textbf{Forma Normal (eixo $OZ$):} $z^2 = \dfrac{x^2}{a^2} + \dfrac{y^2}{b^2}$ \\
    \textbf{Translada em $V(h,j,l)$:} $(z-l)^2 = \dfrac{(x-h)^2}{a^2} + \dfrac{(y-j)^2}{b^2}$ \\
    \textbf{Características:} Vértice $V(0,0,0)$. As secções $z=k$ (com $k \neq 0$) são elipses. As secções $x=k$ ou $y=k$ (com $k \neq 0$) são hipérboles. Se $k=0$, as secções no vértice são retas concorrentes.
  \end{minipage}%
  \hfill
  \begin{minipage}{0.24\textwidth}
    \centering
    \begin{tikzpicture}[scale=0.52]
      \draw[-stealth, gray] (-1.5,0) -- (1.8,0) node[right, black] {\tiny $y$};
      \draw[-stealth, gray] (0,-1.6) -- (0,1.8) node[above, black] {\tiny $z$};
      \draw[-stealth, gray] (0.8,-0.8) -- (-0.8,0.8) node[above left, black] {\tiny $x$};
      
      \draw[thick, primaryColor] (-1.0,1.3) -- (1.0,-1.3);
      \draw[thick, primaryColor] (1.0,1.3) -- (-1.0,-1.3);
      \draw[thick, secondaryColor] (0,1.3) ellipse (1.0 and 0.35);
      \draw[thick, secondaryColor] (0,-1.3) ellipse (1.0 and 0.35);
    \end{tikzpicture}
  \end{minipage}
\end{formulabox}

\begin{formulabox}[top=2pt, bottom=2pt, left=4pt, right=4pt]{6. Hiperbolóide de 1 folha}
  \enlargethispage{1.8cm}
  \begin{minipage}{0.72\textwidth}
    \textbf{Forma Normal (eixo $OZ$):} $\dfrac{x^2}{a^2} + \dfrac{y^2}{b^2} - \dfrac{z^2}{c^2} = 1$ \\
    \textbf{Translada em $C(h,j,l)$:} $\dfrac{(x-h)^2}{a^2} + \dfrac{(y-j)^2}{b^2} - \dfrac{(z-l)^2}{c^2} = 1$ \\
    \textbf{Características:} Eixo de simetria paralelo ao eixo correspondente ao termo negativo ($OZ$ neste caso). 
    As secções $z=k$ são elipses para todo o $k$. As secções $x=k$ ou $y=k$ são hipérboles. É uma superfície regrada.
  \end{minipage}%
  \hfill
  \begin{minipage}{0.24\textwidth}
    \centering
    \begin{tikzpicture}[scale=0.52]
      \draw[-stealth, gray] (-1.5,0) -- (1.8,0) node[right, black] {\tiny $y$};
      \draw[-stealth, gray] (0,-1.6) -- (0,1.8) node[above, black] {\tiny $z$};
      \draw[-stealth, gray] (0.8,-0.8) -- (-0.8,0.8) node[above left, black] {\tiny $x$};
      
      \draw[thick, primaryColor] (-1.0,1.3) to[out=-60, in=120] (-0.6,0) to[out=-120, in=60] (-1.0,-1.3);
      \draw[thick, primaryColor] (1.0,1.3) to[out=-120, in=60] (0.6,0) to[out=-60, in=120] (1.0,-1.3);
      
      \draw[thick, secondaryColor] (0,1.3) ellipse (1.0 and 0.3);
      \draw[thick, secondaryColor] (0,0) ellipse (0.6 and 0.18);
      \draw[thick, secondaryColor] (0,-1.3) ellipse (1.0 and 0.3);
    \end{tikzpicture}
  \end{minipage}
\end{formulabox}

\begin{formulabox}[top=2pt, bottom=2pt, left=4pt, right=4pt]{7. Hiperbolóide de 2 folhas}
  \begin{minipage}{0.72\textwidth}
    \textbf{Forma Normal (eixo $OZ$):} $\dfrac{z^2}{c^2} - \dfrac{x^2}{a^2} - \dfrac{y^2}{b^2} = 1$ \\
    \textbf{Translada em $C(h,j,l)$:} $\dfrac{(z-l)^2}{c^2} - \dfrac{(x-h)^2}{a^2} - \dfrac{(y-j)^2}{b^2} = 1$ \\
    \textbf{Características:} Eixo de simetria paralelo ao eixo do termo positivo ($OZ$ neste caso).
    As secções $x=k$ e $y=k$ são hipérboles. As secções $z=k$ são elipses apenas para $|k| > c$ (não há pontos reais na faixa $-c < z < c$).
  \end{minipage}%
  \hfill
  \begin{minipage}{0.24\textwidth}
    \centering
    \begin{tikzpicture}[scale=0.52]
      \draw[-stealth, gray] (-1.5,0) -- (1.8,0) node[right, black] {\tiny $y$};
      \draw[-stealth, gray] (0,-1.8) -- (0,2.0) node[above, black] {\tiny $z$};
      \draw[-stealth, gray] (0.8,-0.8) -- (-0.8,0.8) node[above left, black] {\tiny $x$};
      
      \draw[thick, primaryColor] (-1.0,1.6) to[out=-60, in=180] (0,0.6) to[out=0, in=-120] (1.0,1.6);
      \draw[thick, secondaryColor] (0,1.6) ellipse (1.0 and 0.3);
      
      \draw[thick, primaryColor] (-1.0,-1.6) to[out=60, in=180] (0,-0.6) to[out=0, in=120] (1.0,-1.6);
      \draw[thick, secondaryColor] (0,-1.6) ellipse (1.0 and 0.3);
    \end{tikzpicture}
  \end{minipage}
\end{formulabox}

\begin{formulabox}[top=2pt, bottom=2pt, left=4pt, right=4pt]{8. Superfície Cilíndrica (Cilindro Elíptico / Circular)}
  \begin{minipage}{0.72\textwidth}
    \textbf{Forma Normal (eixo $OZ$):} $\dfrac{x^2}{a^2} + \dfrac{y^2}{b^2} = r^2$ \\
    \textbf{Translada em $C(h,j)$:} $\dfrac{(x-h)^2}{a^2} + \dfrac{(y-j)^2}{b^2} = r^2$ \\
    \textbf{Características:} Ausência da variável correspondente ao eixo gerador (neste caso, $z$). 
    As secções $z=k$ são elipses idênticas. As secções $x=k$ (com $|k| < ar$) ou $y=k$ (com $|k| < br$) são pares de retas paralelas ao eixo $OZ$.
  \end{minipage}%
  \hfill
  \begin{minipage}{0.24\textwidth}
    \centering
    \begin{tikzpicture}[scale=0.52]
      \draw[-stealth, gray] (-1.5,0) -- (1.8,0) node[right, black] {\tiny $y$};
      \draw[-stealth, gray] (0,-1.6) -- (0,1.8) node[above, black] {\tiny $z$};
      \draw[-stealth, gray] (0.8,-0.8) -- (-0.8,0.8) node[above left, black] {\tiny $x$};
      
      \draw[thick, primaryColor] (-1.0,-1.3) -- (-1.0,1.3);
      \draw[thick, primaryColor] (1.0,-1.3) -- (1.0,1.3);
      \draw[thick, secondaryColor] (0,1.3) ellipse (1.0 and 0.3);
      \draw[thick, secondaryColor] (0,-1.3) ellipse (1.0 and 0.3);
    \end{tikzpicture}
  \end{minipage}
\end{formulabox}

\begin{formulabox}[top=2pt, bottom=2pt, left=4pt, right=4pt]{9. Cilindro Parabólico}
  \begin{minipage}{0.72\textwidth}
    \textbf{Forma Normal (eixo $OZ$):} $y = ax^2$ \\
    \textbf{Translada em $V(h,j)$:} $y-j = a(x-h)^2$ \\
    \textbf{Características:} Ausência de uma das variáveis tridimensionais (neste caso, $z$). A curva diretriz no plano transversal $xOy$ é uma parábola. As secções $z=k$ são parábolas idênticas e as secções $x=k$ ou $y=k$ são retas verticais. \label{box:9}
  \end{minipage}%
  \hfill
  \begin{minipage}{0.24\textwidth}
    \centering
    \begin{tikzpicture}[scale=0.52]
      \draw[-stealth, gray] (-1.5,0) -- (1.8,0) node[right, black] {\tiny $y$};
      \draw[-stealth, gray] (0,-1.6) -- (0,1.8) node[above, black] {\tiny $z$};
      \draw[-stealth, gray] (0.8,-0.8) -- (-0.8,0.8) node[above left, black] {\tiny $x$};
      
      \draw[thick, secondaryColor, domain=-0.9:0.9] plot (\x, {0.8*\x*\x + 0.5});
      \draw[thick, secondaryColor, domain=-0.9:0.9] plot (\x, {0.8*\x*\x - 1.2});
      
      \draw[thick, primaryColor] (-0.9, 1.15) -- (-0.9, -0.55);
      \draw[thick, primaryColor] (0.9, 1.15) -- (0.9, -0.55);
      \draw[thick, primaryColor] (0, 0.5) -- (0, -1.2);
    \end{tikzpicture}
  \end{minipage}
\end{formulabox}


\subsubsection{Guia de Identificação Rápida de Quádricas}
\label{sec:guia}
\begin{enumerate}[label=\textbf{\arabic*.}]
  \item $x^2 + 4y^2 + 9z^2 = 1 \implies$ \textbf{Gráfico VII} (Elipsoide): Todos os coeficientes quadráticos positivos. Semi-eixos $a=1, b=1/2, c=1/3$. Alongado ao longo de $OX$ (eixo do maior semi-eixo).
  \item $9x^2 + 4y^2 + z^2 = 1 \implies$ \textbf{Gráfico IV} (Elipsoide): Todos os coeficientes positivos. Semi-eixos $a=1/3, b=1/2, c=1$. Alongado ao longo do eixo vertical $OZ$.
  \item $x^2 - y^2 + z^2 = 1 \implies$ \textbf{Gráfico II} (Hiperboloide de 1 Folha): Dois termos positivos, um negativo ($-y^2$). O termo negativo indica que o eixo do hiperboloide é paralelo a $OY$.
  \item $-x^2 + y^2 - z^2 = 1 \implies$ \textbf{Gráfico III} (Hiperboloide de 2 Folhas): Apenas um termo positivo ($y^2$), dois negativos. O termo positivo define o eixo de simetria paralelo a $OY$.
  \item $y = 2x^2 + z^2 \implies$ \textbf{Gráfico VI} (Parabolóide Elíptico): Uma variável linear ($y$) e duas quadráticas com o mesmo sinal. Abre-se na direção do eixo positivo $OY$.
  \item $y^2 = x^2 + 2z^2 \implies$ \textbf{Gráfico I} (Cone Elíptico): Equação homogénea ($y^2 - x^2 - 2z^2 = 0$). Vértice na origem e eixo de simetria ao longo de $OY$.
  \item $x^2 + 2z^2 = 1 \implies$ \textbf{Gráfico VIII} (Cilindro Elíptico): Ausência de $y$. Descreve uma elipse no plano $xOz$ projetada infinitamente ao longo do eixo $OY$.
  \item $y = x^2 - z^2 \implies$ \textbf{Gráfico V} (Parabolóide Hiperbólico): Uma variável linear ($y$) e duas quadráticas com sinais opostos. Superfície em forma de sela com ponto de sela na origem.
\end{enumerate}

\newpage

% ==========================================
% Capítulo 3: Curvas e Funções Vetoriais de 1 Parâmetro
% ==========================================
\section{Curvas e Funções Vetoriais de 1 Parâmetro}
\label{chap:curvas_funcoes_vetoriais}

\subsection{Definição e Parametrização}

\noindent\begin{minipage}[t]{0.65\textwidth}
  \begin{definitionbox}{Curvas e Equações Paramétricas}
    Seja $I \subseteq \R$ um intervalo. Um sistema de equações paramétricas descreve uma curva no plano ou no espaço onde as coordenadas são funções de uma variável independente $t \in I$ (parâmetro):
    \begin{itemize}
      \item \textbf{No plano ($\R^2$):} $\begin{cases} x = x(t) \\ y = y(t) \end{cases} \implies C = \{(x(t), y(t)) : t \in I\}$.
      \item \textbf{No espaço ($\R^3$):} $\begin{cases} x = x(t) \\ y = y(t) \\ z = z(t) \end{cases} \implies C = \{(x(t), y(t), z(t)) : t \in I\}$.
      \item \textbf{Orientação:} O sentido do movimento ao longo da curva à medida que $t$ aumenta.
      \item \textbf{Tipos de Curvas (para $I = [a, b]$):}
        \begin{enumerate}
          \item \textbf{Regular:} se $\vec{r}'(t) \neq \vec{0}, \ \forall t \in I$ (sem pontos de paragem/bicos).
          \item \textbf{Suave:} se $\vec{r}(t)$ é de classe $C^1$ em $I$ (diferenciável e com derivada $\vec{r}'(t)$ contínua).
          \item \textbf{Fechada:} se o ponto inicial coincide com o final: $\vec{r}(a) = \vec{r}(b)$.
        \end{enumerate}
    \end{itemize}
  \end{definitionbox}
\end{minipage}%
\hfill
\begin{minipage}[t]{0.33\textwidth}
  \centering
  \vspace{8pt}
  \begin{tikzpicture}[scale=0.9]
    \draw[-stealth, gray] (-2.0,0) -- (2.0,0) node[right, black] {\tiny $x$};
    \draw[-stealth, gray] (0,-1.2) -- (0,1.2) node[above, black] {\tiny $y$};
    
    % Ellipse x^2/2.25 + y^2/0.64 = 1
    \draw[thick, primaryColor] (0,0) ellipse (1.5 and 0.8);
    
    % Orientation arrows (clockwise for x = 5sin(t), y = 2cos(t))
    \draw[primaryColor, -stealth, thick] (0, 0.8) -- (0.75, 0.57);
    \draw[primaryColor, -stealth, thick] (0, -0.8) -- (-0.75, -0.57);
    
    \fill[accentColor] (0, 0.8) circle (1.5pt) node[above right, black] {\tiny $t=0$};
    \fill[accentColor] (1.5, 0) circle (1.5pt) node[below right, black] {\tiny $t=\frac{\pi}{2}$};
    
    \node[font=\tiny, align=center] at (0, -1.4) {Elipse com orientação \\ horária ($x=5\sin t, y=2\cos t$)};
  \end{tikzpicture}
\end{minipage}

\subsection{Geometria Diferencial de Curvas}

\noindent\begin{minipage}[t]{0.65\textwidth}
  \begin{definitionbox}{Cálculo em Funções Vetoriais}
    Seja $\vec{r} : I \subseteq \R \to \R^n$ uma função vetorial.
    \begin{itemize}
      \item \textbf{Limite:} Calcula-se componente a componente: \\
        $\lim_{t \to t_0} \vec{r}(t) = \left( \lim_{t \to t_0} r_1(t), \dots, \lim_{t \to t_0} r_n(t) \right)$.
      \item \textbf{Derivada:} Vetor tangente/velocidade à curva no ponto $t_0$: \\
        $\vec{r}'(t_0) = \lim_{h \to 0} \frac{\vec{r}(t_0+h) - \vec{r}(t_0)}{h} = \left( r_1'(t_0), \dots, r_n'(t_0) \right)$.
    \end{itemize}
  \end{definitionbox}

  \begin{theorembox}{Comprimento de Arco e Reta Tangente}
    Seja $\vec{r}(t)$ uma curva suave em $I = [a, b]$.
    \begin{itemize}
      \item \textbf{Comprimento da Curva ($L$):} Integral do módulo da velocidade: \\
        $L = \int_a^b \|\vec{r}'(t)\| \dd t = \int_a^b \sqrt{(r_1'(t))^2 + \dots + (r_n'(t))^2} \dd t$.
      \item \textbf{Reta Tangente:} Reta tangente à curva no ponto $P_0 = \vec{r}(t_0)$ na direção de $\vec{r}'(t_0) \neq \vec{0}$, parametrizada por: \\
        $X(\lambda) = P_0 + \lambda \vec{r}'(t_0), \quad \lambda \in \R$.
    \end{itemize}
  \end{theorembox}
\end{minipage}%
\hfill
\begin{minipage}[t]{0.33\textwidth}
  \centering
  \vspace{8pt}
  \begin{tikzpicture}[scale=0.9, x={(-0.4cm,-0.4cm)}, y={(1.0cm,0cm)}, z={(0cm,1.0cm)}]
    % Axes
    \draw[-stealth, gray] (0,0,0) -- (1.5,0,0) node[below left, black] {\tiny $x$};
    \draw[-stealth, gray] (0,0,0) -- (0,1.8,0) node[right, black] {\tiny $y$};
    \draw[-stealth, gray] (0,0,0) -- (0,0,3.3) node[above, black] {\tiny $z$};
    
    % Helix curve r(t) = (cos(t), sin(t), 0.35*t)
    \draw[thick, primaryColor, variable=\t, domain=0:2.5*pi, samples=80]
      plot ({cos(\t r)}, {sin(\t r)}, {0.35*\t});
    
    % Arrow showing orientation
    \draw[primaryColor, -stealth, thick, variable=\t, domain=1.0*pi:1.15*pi, samples=10]
      plot ({cos(\t r)}, {sin(\t r)}, {0.35*\t});
    
    \fill[accentColor] ({cos(0)}, {sin(0)}, 0) circle (1.5pt) node[right, black] {\tiny $t=0$};
    \fill[accentColor] ({cos(2.5*pi r)}, {sin(2.5*pi r)}, {0.35*2.5*pi}) circle (1.5pt) node[right, black] {\tiny $t=2.5\pi$};
    
    \node[font=\tiny, align=center] at (0,0,3.9) {Hélice Circular: \\ $x=3\sin t$, $y=3\cos t$ \\ $z=0.5t$, $t\in[0,2\pi]$};
  \end{tikzpicture}
\end{minipage}


% ==========================================
% Capítulo 4: Limites e Continuidade em $\R^n$
% ==========================================
\section{Limites e Continuidade em \texorpdfstring{$\R^n$}{R\textasciicircum n}}
\label{chap:limites_continuidade}

\subsection{Topologia e Sucessões em \texorpdfstring{$\R^n$}{R\textasciicircum n}}

\begin{minipage}[t]{0.65\textwidth}
  \begin{definitionbox}{Conceitos Topológicos em $\R^n$}
    Seja $A \subseteq \R^n$ um conjunto e $a \in \R^n$.
    \begin{itemize}
      \item \textbf{Bola Aberta de centro $a$ e raio $\delta > 0$:} $B_\delta(a) = \{x \in \R^n : \|x - a\| < \delta\}$.
      \item \textbf{Ponto Aderente e Fecho ($\bar{A}$):} $a$ é aderente a $A$ se $\forall \delta > 0, B_\delta(a) \cap A \neq \emptyset$. O fecho $\bar{A}$ é o conjunto de todos os pontos aderentes a $A$.
      \item \textbf{Ponto Interior e Interior ($\text{int}(A)$):} $a$ é interior a $A$ se $\exists \delta > 0, B_\delta(a) \subseteq A$. O interior $\text{int}(A)$ é o conjunto de todos os pontos interiores de $A$.
      \item \textbf{Ponto de Fronteira e Fronteira ($\partial A$):} $a$ é de fronteira se não for interior nem exterior: $\forall \delta > 0$, $B_\delta(a) \cap A \neq \emptyset \wedge B_\delta(a) \cap (\R^n \setminus A) \neq \emptyset$.
      \item \textbf{Ponto Exterior e Exterior:} $a$ é exterior se for interior ao complementar de $A$: $\exists \delta > 0, B_\delta(a) \subseteq \R^n \setminus A$.
      \item \textbf{Conjunto Aberto e Fechado:}
        \begin{itemize}
          \item $A$ é \textbf{aberto} se $A = \text{int}(A)$ (todos os seus pontos são interiores).
          \item $A$ é \textbf{fechado} se $A = \bar{A}$ (contém todos os seus pontos aderentes).
        \end{itemize}
      \item \textbf{Conjunto Compacto:} $A$ é compacto se for \textbf{fechado e limitado}.
    \end{itemize}
  \end{definitionbox}
\end{minipage}%
\hfill
\begin{minipage}[t]{0.33\textwidth}
  \centering
  \vspace{8pt}
  \begin{tikzpicture}[scale=1.0]
    % Shade region A (light teal)
    \fill[secondaryColor!15] (0,0) circle (1.2);
    % Boundary of A (dashed line)
    \draw[thick, dashed, secondaryColor] (0,0) circle (1.2) node[above right=0.9cm] {\small $A$};
    
    % Interior point
    \coordinate (Int) at (-0.3, 0.2);
    \fill[primaryColor] (Int) circle (1.5pt) node[below] {\tiny $x_{\text{int}}$};
    \draw[dashed, primaryColor] (Int) circle (0.3);
    
    % Boundary point
    \coordinate (Bnd) at (0.97, 0.7); % on the circle x^2+y^2 = 1.2^2 approximately
    \fill[accentColor] (Bnd) circle (1.5pt) node[right] {\tiny $x_{\text{fr}}$};
    \draw[dashed, accentColor] (Bnd) circle (0.3);
    
    % Exterior point
    \coordinate (Ext) at (1.7, -0.5);
    \fill[darkGray] (Ext) circle (1.5pt) node[below] {\tiny $x_{\text{ext}}$};
    \draw[dashed, darkGray] (Ext) circle (0.3);
    
    % Labels
    \node[secondaryColor] at (0, -0.4) {\tiny $\text{int}(A)$};
    \node[accentColor] at (1.3, 1.1) {\tiny $\partial A$};
  \end{tikzpicture}
  
  \begin{theorembox}{Sucessões em $\R^n$}
    Uma sucessão em $\R^n$ é $u : \N \to \R^n$, $m \mapsto u_m = (u_m^1, \dots, u_m^n)$.
    \begin{itemize}
      \item \textbf{Exemplo:} $u_m = (m^2, 1/m)$.
      \item \textbf{Limite:} $\lim u_m = a \iff \forall \epsilon > 0, \exists p \in \N, \forall m \ge p \implies \|u_m - a\| < \epsilon$.
      \item \textbf{Convergência:} $\lim u_m = a \iff \lim u_m^i = a_i, \ \forall i \in \{1,\dots,n\}$.
    \end{itemize}
  \end{theorembox}
\end{minipage}

\subsection{Funções de Várias Variáveis e Domínios}

\begin{minipage}[t]{0.65\textwidth}
  \begin{definitionbox}{Domínio, Gráfico e Conjuntos de Nível}
    Seja $f: D \subseteq \R^n \to \R$ uma função real de várias variáveis.
    \begin{itemize}
      \item \textbf{Domínio ($D$):} O maior subconjunto de $\R^n$ onde a expressão $f(x_1, \dots, x_n)$ tem significado matemático.
      \item \textbf{Contra-domínio ($f(D)$) e Gráfico ($G$):} \\
        $f(D) = \{z \in \R : z = f(x) \wedge x \in D\}$ \\
        $G = \{(x_1, \dots, x_n, z) \in \R^{n+1} : z = f(x) \wedge x \in D\}$
      \item \textbf{Conjunto de Nível ($C_k$):} $C_k = \{x \in \R^n : f(x) = k\}$ para $k \in f(D)$. Em $\R^2$, é a \textbf{curva de nível}; em $\R^3$, é a \textbf{superfície de nível}.
    \end{itemize}
  \end{definitionbox}

  \begin{theorembox}{Cuidados no Cálculo de Domínios}
    Garantir as seguintes restrições em todo o domínio $D$:
    \begin{enumerate}
      \item \textbf{Denominadores} não nulos: $P(x)/Q(x) \implies Q(x) \neq 0$.
      \item \textbf{Raízes de índice par} não negativas: $\sqrt[2k]{g(x)} \implies g(x) \geq 0$.
      \item \textbf{Logaritmos} com argumento positivo: $\ln(g(x)) \implies g(x) > 0$.
      \item \textbf{Arcosseno e Arcocosseno} em $[-1, 1]$: $-1 \leq g(x) \leq 1$.
      \item \textbf{Tangente:} $g(x) \neq \pi/2 + k\pi, \ k \in \Z$.
    \end{enumerate}
  \end{theorembox}
\end{minipage}%
\hfill
\begin{minipage}[t]{0.33\textwidth}
  \centering
  \vspace{8pt}
  \begin{tikzpicture}[scale=1.0]
    \draw[-stealth, gray] (-2.0,0) -- (2.0,0) node[right, black] {\tiny $x$};
    \draw[-stealth, gray] (0,-1.5) -- (0,1.5) node[above, black] {\tiny $y$};
    
    % Concentric ellipses for f(x,y) = x^2 + 2y^2 = k
    \draw[thick, primaryColor] (0,0) ellipse (0.7 and 0.49);
    \node[primaryColor, above right] at (0.49, 0.35) {\tiny $C_{k_1}$};
    
    \draw[thick, secondaryColor] (0,0) ellipse (1.2 and 0.84);
    \node[secondaryColor, above right] at (0.84, 0.59) {\tiny $C_{k_2}$};
    
    \draw[thick, accentColor] (0,0) ellipse (1.7 and 1.19);
    \node[accentColor, above right] at (1.19, 0.83) {\tiny $C_{k_3}$};
    
    \node[draw, fill=white, font=\tiny] at (0,1.3) {$f(x,y)=k$ ($k_1 < k_2 < k_3$)};
  \end{tikzpicture}
  \vspace{4pt}
  \begin{center}
    \small \textbf{Curvas de Nível ($C_k$):} Secções horizontais do gráfico projetadas no plano $xOy$.
  \end{center}
\end{minipage}

\subsection{Limites Multivariáveis}

\begin{definitionbox}{Limite de Funções Reais e Vetoriais}
  Seja $f : D \subseteq \R^n \to \R^p$ uma função e $a \in \bar{D}$.
  \begin{itemize}
    \item \textbf{Limite segundo Cauchy:} $\lim_{x \to a} f(x) = b \iff \forall \delta > 0, \exists \epsilon > 0, x \in D \cap B_\epsilon(a) \implies f(x) \in B_\delta(b)$, i.e., $\forall \delta > 0, \exists \epsilon > 0, x \in D \wedge \|x - a\| < \epsilon \implies \|f(x) - b\| < \delta$.
    \item \textbf{Limite segundo Heine:} $\lim_{x\to a} f(x) = b \iff \forall (x_m) \subset D$ com $x_m \to a$, tem-se $f(x_m) \to b$.
    \item \textbf{Convergência por Componentes:} $\lim_{x \to a} f(x) = b \iff \lim_{x \to a} f_i(x) = b_i, \ \forall i \in \{1, \dots, p\}$.
  \end{itemize}
\end{definitionbox}

\begin{theorembox}{Propriedades Operatórias e Squeeze Theorem}
  Sejam $f, g, h : D \subseteq \R^n \to \R$ com limites finitos em $a \in \bar{D}$.
  \begin{itemize}
    \item \textbf{Álgebra de Limites:} $\lim(f \pm g) = \lim f \pm \lim g$; $\lim(f \cdot g) = \lim f \cdot \lim g$; $\lim(f/g) = \lim f / \lim g$ (se $\lim g \neq 0$).
    \item \textbf{Carácter Local:} Se $f(x) = g(x)$ em $B_\epsilon(a) \cap D \implies \lim f = \lim g$.
    \item \textbf{Enquadramento (Squeeze):} Se $f(x) \leq h(x) \leq g(x)$ em $B_\epsilon(a) \cap D$ e $\lim f = \lim g = b \implies \lim h = b$.
    \item \textbf{Corolário (Infinitésimo por Limitada):} Se $\lim g = 0$ e $f$ é limitada ($|f| \leq L$) em $B_\epsilon(a) \cap D \implies \lim(f \cdot g) = 0$.
    \item \textbf{União de Domínios:} Se $D = \cup D_i$ e $\lim_{x\to a, x\in D_i} f(x) = b \ \forall i \implies \lim_{x\to a} f(x) = b$.
  \end{itemize}
\end{theorembox}

\begin{minipage}[t]{0.65\textwidth}
  \begin{definitionbox}{Limites Relativos e Direcionais}
    Seja $f : D \subseteq \R^n \to \R$ e $A \subset D$ com $a \in \bar{A}$.
    \begin{itemize}
      \item \textbf{Limite Relativo a $A$:} $\lim_{\substack{x \to a \\ x \in A}} f(x) = b \iff$ a restrição de $f$ a $A$ tem limite $b$. Pelo critério de Heine, $\forall u_m \in A, u_m \to a \implies f(u_m) \to b$.
      \item \textbf{Limites Direcionais (em $\R^2$):} Aproximação a $a(a_1, a_2)$ por retas:
        $A_m = \{(x, y) \in \R^2 : y - a_2 = m(x - a_1)\}$ (declive $m$) ou a reta vertical $A_\infty = \{(x, y) \in \R^2 : x = a_1\}$.
      \item \textbf{Caminhos Parabólicos:} $A_k = \{(x, y) \in \R^2 : y - a_2 = k(x - a_1)^2\}$.
    \end{itemize}
  \end{definitionbox}

  \begin{theorembox}{Teoremas de Não-Existência de Limite}
    Sejam $A_1, A_2 \subseteq D$ com $a \in \bar{A}_1 \cap \bar{A}_2$.
    \begin{enumerate}
      \item Se os limites relativos diferem ou não existem: $\lim_{\substack{x \to a \\ x \in A_1}} f(x) \neq \lim_{\substack{x \to a \\ x \in A_2}} f(x) \implies \not\exists \lim_{x\to a} f(x)$.
      \item Se o limite direcional $\lim_{\substack{(x,y)\to a \\ (x,y)\in A_m}} f(x,y)$ depender de $m \implies \not\exists \lim f$.
    \end{enumerate}
  \end{theorembox}
\end{minipage}%
\hfill
\begin{minipage}[t]{0.33\textwidth}
  \centering
  \vspace{8pt}
  \begin{tikzpicture}[scale=1.0]
    \draw[-stealth, gray] (-1.8,0) -- (1.8,0) node[right, black] {\tiny $x$};
    \draw[-stealth, gray] (0,-1.8) -- (0,1.8) node[above, black] {\tiny $y$};
    
    % Point a at origin
    \fill[accentColor] (0,0) circle (1.8pt) node[below right, black] {\tiny $a(0,0)$};
    
    % Path y = mx (lines)
    \draw[domain=-1.6:1.6, smooth, variable=\x, primaryColor, thick] plot (\x, {0.5*\x}) node[right] {\tiny $y=mx$};
    \draw[domain=-1.3:1.3, smooth, variable=\x, primaryColor, thick] plot (\x, {-\x}) node[left] {\tiny $y=-x$};
    
    % Path y = kx^2 (parabola)
    \draw[domain=-1.2:1.2, smooth, variable=\x, secondaryColor, thick] plot (\x, {\x*\x}) node[above] {\tiny $y=kx^2$};
    
    % Arrows along the paths to show approach
    \draw[primaryColor, -stealth] (1.1, 0.55) -- (0.6, 0.3);
    \draw[primaryColor, -stealth] (-0.9, 0.9) -- (-0.5, 0.5);
    \draw[secondaryColor, -stealth] (0.9, 0.81) -- (0.5, 0.25);
  \end{tikzpicture}
  \vspace{2pt}
  \begin{center}
    \small \textbf{Caminhos de Aproximação:} O limite deve ser independente do caminho para existir.
  \end{center}
\end{minipage}

\begin{formulabox}{Estratégias Práticas para Resolução de Limites}
  Para $\lim_{(x,y) \to (a_1, a_2)} f(x,y) = b$:
  \begin{itemize}
    \item \textbf{1. Testar Não-Existência:} Analisar limites direcionais ou iterados. Se dependerem de $m$ ou forem diferentes, não existe.
    \item \textbf{2. Enquadramento (Infinitésimos):} Se candidato a $b$ existir, majorar $|f(x,y) - b| \leq h(x,y)$ com $\lim h = 0$. Utilizar $|x| \leq \sqrt{x^2+y^2}$ e $|y| \leq \sqrt{x^2+y^2}$.
    \item \textbf{3. Coordenadas Polares (em $(0,0)$):} Substituir $x = r\cos\theta, y = r\sin\theta$. O limite é $b \iff |f(r\cos\theta, r\sin\theta) - b| \leq h(r)$ com $\lim_{r \to 0^+} h(r) = 0$ independentemente de $\theta$.
  \end{itemize}
\end{formulabox}

\subsection{Continuidade e Teoremas Fundamentais}

\begin{definitionbox}{Continuidade e Prolongamento}
  Seja $f : D \subseteq \R^n \to \R^p$ uma função e $a \in D$.
  \begin{itemize}
    \item \textbf{Continuidade em um Ponto:} $f$ é contínua em $a$ se $\lim_{x\to a} f(x) = f(a)$.
    \item \textbf{Continuidade em um Conjunto:} $f$ é contínua em $A$ se for contínua em todos os pontos de $A$.
    \item \textbf{Prolongamento por Continuidade:} $a \in \bar{D} \setminus D$. Se $\lim_{x\to a} f(x) = b$, então $\bar{f}(x) = f(x)$ para $x \in D$ e $\bar{f}(a) = b$ é contínua.
    \item \textbf{Descontinuidade Removível:} $a \in D$ mas $\lim_{x\to a} f(x) = b \neq f(a)$. Basta redefinir $f(a) = b$.
  \end{itemize}
\end{definitionbox}

\begin{theorembox}{Álgebra e Composição de Contínuas}
  \begin{itemize}
    \item \textbf{Operações:} Se $f, g$ contínuas em $a \implies f \pm g, f \cdot g, f/g$ (se $g(a) \neq 0$) são contínuas em $a$.
    \item \textbf{Composição:} $g$ contínua em $a$ e $f$ contínua em $g(a) \implies f \circ g$ é contínua em $a$.
    \item \textbf{Funções Elementares:} Constantes, projeções, polinómios, exponenciais, logaritmos e trigonométricas são contínuas nos seus domínios.
  \end{itemize}
\end{theorembox}

\begin{theorembox}{Teoremas Fundamentais das Funções Contínuas}
  \begin{itemize}
    \item \textbf{Preservação da Compacidade:} $K$ compacto e $f$ contínua $\implies f(K)$ compacto.
    \item \textbf{Teorema de Weierstrass (Extremos Absolutos):} $K$ compacto e $f$ contínua $\implies \exists x_1, x_2 \in K$ tais que $f(x_1) \leq f(x) \leq f(x_2), \ \forall x \in K$.
    \item \textbf{Teorema de Bolzano (Valor Intermédio):} $f$ contínua em $D$, $a, b \in D$, $[a, b] \subseteq D \implies \forall k \in [f(a), f(b)], \exists c \in [a, b]$ tal que $f(c) = k$.
  \end{itemize}
\end{theorembox}


% ==========================================
% Capítulo 5: Cálculo Diferencial em $\R^n$
% ==========================================
\section{Cálculo Diferencial em \texorpdfstring{$\R^n$}{R\textasciicircum n}}
\label{chap:calculo_diferencial}

\subsection{Derivadas Parciais e Estruturas de Derivação}

\noindent\begin{definitionbox}{Derivadas Parciais e de Ordem Superior}
  Seja $f : D \subseteq \R^n \to \R$ e $a = (a_1, \dots, a_n) \in \text{int}(D)$.
  \begin{itemize}
    \item \textbf{Derivada Parcial em ordem a $x_i$:} Taxa de variação em relação a $x_i$ no ponto $a$:
      \[ \frac{\partial f}{\partial x_i}(a) = \lim_{h\to 0} \frac{f(a_1, \dots, a_i+h, \dots, a_n) - f(a_1, \dots, a_n)}{h} \]
    \item \textbf{Derivadas de Ordem Superior:} Derivação sucessiva de funções derivadas parciais:
      \[ \frac{\partial^2 f}{\partial x_j \partial x_i}(a) = \frac{\partial}{\partial x_j} \left( \frac{\partial f}{\partial x_i} \right) (a), \quad \frac{\partial^2 f}{\partial x_i^2}(a) = \frac{\partial}{\partial x_i} \left( \frac{\partial f}{\partial x_i} \right) (a) \]
    \item \textbf{Classe $C^p$:} $f \in C^p(A)$ num aberto $A$ se as derivadas parciais de $f$ até à ordem $p$ existem e são contínuas em $A$. $f \in C^\infty(A) \iff f \in C^p(A), \forall p \in \N$.
  \end{itemize}
\end{definitionbox}

\noindent\begin{theorembox}{Teorema de Schwarz, Gradiente e Hessiana}
  \begin{itemize}
    \item \textbf{Teorema de Schwarz (Simetria das Mistas):} Se $f \in C^2(A)$, a ordem das derivações parciais de segunda ordem pode ser trocada:
      \[ \frac{\partial^2 f}{\partial x_j \partial x_i}(a) = \frac{\partial^2 f}{\partial x_i \partial x_j}(a) \]
    \item \textbf{Vetor Gradiente ($\nabla f$):} Vetor coluna das derivadas parciais de primeira ordem:
      \[ \nabla f(x) = \left[ \frac{\partial f}{\partial x_1}(x), \dots, \frac{\partial f}{\partial x_n}(x) \right]^T \]
    \item \textbf{Matriz Hessiana ($\text{Hess}\,f$):} Matriz quadrada das derivadas parciais de 2.ª ordem:
      \[ \text{Hess}\,f(x) = \left[ \frac{\partial^2 f}{\partial x_i \partial x_j}(x) \right]_{n \times n} \]
      Pelo Teorema de Schwarz, se $f \in C^2(A)$, a matriz Hessiana é \textbf{simétrica}.
  \end{itemize}
\end{theorembox}

\subsection{Diferenciabilidade e Diferencial}

\noindent\begin{definitionbox}{Diferenciabilidade em $\R^n$}
  Seja $f : D \subseteq \R^n \to \R$ e $a \in \text{int}(D)$.
  \begin{itemize}
    \item \textbf{Definição:} $f$ é diferenciável em $a$ se existe aplicação linear $L_a : \R^n \to \R$ tal que:
      \[ f(a+h) = f(a) + L_a(h) + \|h\|\epsilon(h), \quad \text{com } \lim_{h\to 0}\epsilon(h) = 0 \]
    \item \textbf{Diferenciabilidade $\implies$ Continuidade:} Se $f$ é diferenciável em $a$, então $f$ é contínua em $a$. (A recíproca é falsa).
    \item \textbf{Propriedade:} Se $f$ é diferenciável em $a$, todas as parciais existem e:
      \[ \dd f(a)(h) = L_a(h) = \nabla f(a)^T h = \frac{\partial f}{\partial x_1}(a)h_1 + \dots + \frac{\partial f}{\partial x_n}(a)h_n \]
  \end{itemize}
\end{definitionbox}

\noindent\begin{theorembox}{Condições de Diferenciabilidade}
  \begin{itemize}
    \item \textbf{Em $\R$ vs $\R^n$:} Em $\R$, diferenciável $\iff$ existe derivada. Em $\R^n$, a existência de todas as derivadas parciais \textbf{não garante} diferenciabilidade (nem continuidade).
    \item \textbf{Condição Suficiente (Classe $C^1$):} Se as parciais $\frac{\partial f}{\partial x_i}$ existem e são \textbf{contínuas} em $A$ aberto $\implies f$ é diferenciável em todos os pontos de $A$.
  \end{itemize}
\end{theorembox}

\noindent\begin{formulabox}{Guia Prático: Como Provar a Diferenciabilidade num Ponto}
  Para verificar se $f(x)$ é diferenciável no ponto $a \in \text{int}(D)$:
  \begin{enumerate}
    \item \textbf{Passo 1:} Calcular $\frac{\partial f}{\partial x_i}(a)$ para todo $i \in \{1,\dots,n\}$. Se alguma parcial não existir, $f$ \textbf{não é diferenciável} em $a$.
    \item \textbf{Passo 2:} Obter o candidato a diferencial: $\dd f(a)(h) = \nabla f(a)^T h$.
    \item \textbf{Passo 3:} Calcular o limite do erro $\epsilon(h)$ quando $h \to 0$:
      \[ \lim_{h \to \vec{0}} \frac{f(a+h) - f(a) - \dd f(a)(h)}{\|h\|} = 0 \]
      Se o limite existir e for igual a 0, $f$ é diferenciável em $a$; caso contrário, não é.
  \end{enumerate}
\end{formulabox}

\subsection{Derivadas segundo um Vetor e Direcionais}

\noindent\begin{definitionbox}{Derivada segundo um Vetor e Direcional}
  Seja $f : D \subseteq \R^n \to \R$, $a \in \text{int}(D)$ e $\vec{u} \in \R^n$.
  \begin{itemize}
    \item \textbf{Derivada segundo um vetor $\vec{u}$:} O limite de variação média direcional:
      \[ f'_{\vec{u}}(a) = \lim_{t \to 0} \frac{f(a + t\vec{u}) - f(a)}{t} \]
    \item \textbf{Derivada Direcional:} Se $\vec{u}$ for unitário ($\|\vec{u}\| = 1$), a derivada designa-se por derivada direcional de $f$ em $a$ na direção de $\vec{u}$. Representa a taxa de variação na direção de $\vec{u}$.
    \item \textbf{Propriedade:} A derivada segundo os vetores da base padrão $\vec{e}_i$ corresponde às derivadas parciais: $f'_{\vec{e}_i}(a) = \frac{\partial f}{\partial x_i}(a)$.
  \end{itemize}
\end{definitionbox}

\noindent\begin{theorembox}{Diferenciabilidade, Direções de Máximo e Mínimo}
  \begin{itemize}
    \item \textbf{Teorema:} Se $f$ é diferenciável em $a$, então admite derivada segundo qualquer vetor $\vec{u}$ e esta calcula-se diretamente pelo produto escalar:
      \[ f'_{\vec{u}}(a) = \dd f(a)(\vec{u}) = \nabla f(a)^T \vec{u} = \nabla f(a) \cdot \vec{u} \]
      \textit{Nota:} Uma função pode admitir derivada segundo qualquer vetor num ponto e \textbf{não ser diferenciável} (ou sequer contínua) nesse ponto.
    \item \textbf{Crescimento Máximo (Taxa Máxima):} Ocorre na direção e sentido do vetor gradiente $\nabla f(a)$ (se $\nabla f(a) \neq \vec{0}$). O valor máximo da derivada direcional é $\|\nabla f(a)\|$.
    \item \textbf{Decrescimento Máximo (Taxa Mínima):} Ocorre na direção e sentido oposto ao gradiente $-\nabla f(a)$. O valor mínimo da derivada direcional é $-\|\nabla f(a)\|$.
  \end{itemize}
\end{theorembox}

\subsection{Funções Vetoriais, Plano Tangente e Aproximações}

\noindent\begin{definitionbox}{Diferenciabilidade de Funções Vetoriais}
  Seja $f : D \subseteq \R^n \to \R^p$ com componentes $f(x) = (f_1(x), \dots, f_p(x))$.
  \begin{itemize}
    \item \textbf{Diferenciabilidade Componente a Componente:} $f$ é diferenciável em $a \in \text{int}(D)$ se e só se cada $f_i$ for diferenciável em $a$ para todo o $i \in \{1,\dots,p\}$.
    \item \textbf{Matriz Jacobiana ($\text{Jac}\,f$):} Matriz $p \times n$ das derivadas parciais das componentes:
      \[ \text{Jac}\,f(a) = \begin{bmatrix} \frac{\partial f_1}{\partial x_1}(a) & \dots & \frac{\partial f_1}{\partial x_n}(a) \\ \vdots & \ddots & \vdots \\ \frac{\partial f_p}{\partial x_1}(a) & \dots & \frac{\partial f_p}{\partial x_n}(a) \end{bmatrix} = \begin{bmatrix} \nabla f_1(a)^T \\ \vdots \\ \nabla f_p(a)^T \end{bmatrix} \]
    \item \textbf{Diferencial Vetorial:} Calculado matricialmente por $\dd f(a)(h) = \text{Jac}\,f(a)h$.
  \end{itemize}
\end{definitionbox}

\noindent\begin{theorembox}{Regra da Cadeia e Operações com Diferenciais}
  Sejam $f, g : \R^n \to \R^p$ diferenciáveis em $a \in \text{int}(D)$.
  \begin{itemize}
    \item \textbf{Álgebra de Diferenciais:} $\dd(f \pm g)(a) = \dd f(a) \pm \dd g(a)$. Para $p=1$:
      \[ \dd(f \cdot g)(a) = g(a)\dd f(a) + f(a)\dd g(a), \quad \dd(f/g)(a) = \frac{g(a)\dd f(a) - f(a)\dd g(a)}{g(a)^2} \ (\text{se } g(a)\neq 0) \]
    \item \textbf{Regra da Cadeia (Composição):} Se $g : \R^n \to \R^p$ é diferenciável em $a$ e $f : \R^p \to \R^q$ é diferenciável em $g(a) \implies f \circ g : \R^n \to \R^q$ é diferenciável em $a$:
      \[ \dd(f \circ g)(a) = \dd f(g(a)) \circ \dd g(a) \implies \text{Jac}(f \circ g)(a) = \text{Jac}\,f(g(a)) \times \text{Jac}\,g(a) \]
  \end{itemize}
\end{theorembox}

\noindent\begin{formulabox}{Aproximação Linear e Plano Tangente}
  Seja $f : D \subseteq \R^2 \to \R$ diferenciável em $(x_0, y_0) \in \text{int}(D)$.
  \begin{itemize}
    \item \textbf{Aproximação Linear ($L$):} Para pontos $(x,y)$ próximos de $(x_0, y_0)$, $f(x, y) \approx L(x, y)$:
      \[ L(x, y) = f(x_0, y_0) + \frac{\partial f}{\partial x}(x_0, y_0)(x - x_0) + \frac{\partial f}{\partial y}(x_0, y_0)(y - y_0) \]
      O erro $E(x,y) = f(x,y) - L(x,y)$ satisfaz: $\lim_{(x,y)\to(x_0,y_0)} \frac{E(x,y)}{\|(x,y) - (x_0,y_0)\|} = 0$.
    \item \textbf{Plano Tangente:} A equação do plano tangente ao gráfico de $f$ em $(x_0, y_0, f(x_0,y_0))$ é:
      \[ z = f(x_0, y_0) + \frac{\partial f}{\partial x}(x_0, y_0)(x - x_0) + \frac{\partial f}{\partial y}(x_0, y_0)(y - y_0) \]
    \item \textbf{Variação Aproximada:} $\dd f(x_0, y_0)(x - x_0, y - y_0) \approx f(x, y) - f(x_0, y_0)$.
  \end{itemize}
\end{formulabox}


% ==========================================
% Capítulo 6: Extremos Relativos e Condicionados
% ==========================================
\section{Extremos Relativos e Condicionados}
\label{chap:extremos}

\subsection{Extremos Livres}
Pontos críticos, matriz Hessiana e classificação de pontos (máximo, mínimo ou sela).

\subsection{Extremos Condicionados (Multiplicadores de Lagrange)}
Otimização com restrições de igualdade utilizando multiplicadores de Lagrange.


% ==========================================
% Capítulo 7: Teorema da Função Implícita e Inversa
% ==========================================
\section{Teorema da Função Implícita e Inversa}
\label{chap:funcao_implicita_inversa}

\subsection{Teorema da Função Inversa}
Condições para a invertibilidade local de funções de $\R^n \to \R^n$ e cálculo da derivada da inversa.

\subsection{Teorema da Função Implícita}
Existência e diferenciação de funções definidas implicitamente por sistemas de equações.


% ==========================================
% Capítulo 8: Integrais Duplos
% ==========================================
\section{Integrais Duplos}
\label{chap:integrais_duplos}

\subsection{Integração em Retângulos e Regiões Gerais}
Definição de integral duplo e Teorema de Fubini.

\subsection{Mudança de Variáveis (Coordenadas Polares)}
Transformação de coordenadas e aplicação do Jacobiano em integrais duplos.


% ==========================================
% Capítulo 9: Integrais Triplos
% ==========================================
\section{Integrais Triplos}
\label{chap:integrais_triplos}

\subsection{Integração em Regiões Tridimensionais}
Definição, ordem de integração e cálculo de volumes.

\subsection{Mudanças de Variáveis}
Coordenadas cilíndricas e esféricas com os respetivos Jacobianos.

\end{document}
